\section{Principal divisors and degree}\label{sec:main}
% ===============================================================

We now prove Theorem~\ref{thmA} in two halves.

\subsection{Every principal divisor has degree zero}

\begin{smjproposition}\label{prop:degprin}
For every $f \in \CC(z)^\times$ we have $\deg(\divv(f)) = 0$. Equivalently,
$\Prin(\PP) \subseteq \ker(\deg)$.
\end{smjproposition}

\begin{proof}
By Proposition~\ref{prop:reduced} write $f$ in reduced form,
\[
f(z) = c\, \frac{\prod_{i=1}^{r}(z-a_i)^{m_i}}{\prod_{j=1}^{s}(z-b_j)^{n_j}},
\qquad c \in \CC^\times,
\]
with all $r+s$ points pairwise distinct. By the last assertion of
Proposition~\ref{prop:reduced}, the zeros of $f$ in $\CC$ are exactly the $a_i$
with $\ord_{a_i}(f) = m_i$, the poles in $\CC$ are exactly the $b_j$ with
$\ord_{b_j}(f) = -n_j$, and $\ord_p(f) = 0$ at all other finite $p$. Set
$M := \sum_i m_i$ and $N := \sum_j n_j$.

It remains to compute $\ord_\infty(f)$. Substituting $z = 1/w$,
\[
f(1/w)
= c\,\frac{\prod_i (w^{-1}-a_i)^{m_i}}{\prod_j (w^{-1}-b_j)^{n_j}}
= c\,\frac{\prod_i w^{-m_i}(1-a_iw)^{m_i}}{\prod_j w^{-n_j}(1-b_jw)^{n_j}}
= w^{\,N-M}\, g(w),
\]
where
\[
g(w) := c\,\frac{\prod_i (1-a_iw)^{m_i}}{\prod_j (1-b_jw)^{n_j}} .
\]
Each factor $1 - a_i w$ and $1 - b_j w$ takes the value $1$ at $w = 0$, so $g$ is
holomorphic near $w = 0$ with $g(0) = c \neq 0$. Thus $f(1/w) = w^{N-M}g(w)$ is a
local factorisation at $w = 0$ in the sense of Proposition~\ref{prop:ord}, and
therefore
\[
\ord_\infty(f) = N - M = \sum_{j=1}^{s} n_j - \sum_{i=1}^{r} m_i .
\]
Adding up all the local contributions,
\[
\deg(\divv(f)) = \sum_{i=1}^{r} m_i - \sum_{j=1}^{s} n_j + \ord_\infty(f)
= M - N + (N - M) = 0 . \qedhere
\]
\end{proof}

\begin{remark}
This is the precise sense in which a meromorphic function on $\PP$ has ``as many
zeros as poles''. Specialising to a polynomial $P$ recovers the Fundamental
Theorem of Algebra in the form of Example~\ref{ex:poly}: the $\deg P$ roots in
$\CC$ are balanced by a pole of order $\deg P$ at infinity.
\end{remark}

\subsection{Every divisor of degree zero is principal}

\begin{smjproposition}\label{prop:degzero}
Let $D \in \Div(\PP)$ with $\deg(D) = 0$. Then $D$ is principal.
\end{smjproposition}

\begin{proof}
Write $D = \sum_{k=1}^{r} n_k (a_k) + n_\infty(\infty)$ with
$a_1, \dots, a_r \in \CC$ distinct and $n_k, n_\infty \in \ZZ$; every divisor has
this shape, the point $\infty$ being separated off from the finite ones. Since
$\deg(D) = \sum_k n_k + n_\infty = 0$ we have $n_\infty = -\sum_k n_k$, and so
\[
D \;=\; \sum_{k=1}^{r} n_k (a_k) - \Big(\sum_{k=1}^{r} n_k\Big)(\infty)
\;=\; \sum_{k=1}^{r} n_k\big[(a_k)-(\infty)\big].
\]
By Proposition~\ref{prop:divlinear} each bracket is $\divv(z-a_k)$, and $\divv$
is a homomorphism (Proposition~\ref{prop:divhom}), so
\[
D \;=\; \sum_{k=1}^{r} n_k\, \divv(z-a_k)
\;=\; \divv\Big(\prod_{k=1}^{r}(z-a_k)^{n_k}\Big),
\]
where negative exponents are read as factors in the denominator. Hence $D$ is
principal.
\end{proof}

\begin{remark}
The proof is constructive: given a degree-zero divisor it writes down a function
realising it. Note also that the argument handles the coefficient at $\infty$
automatically, by using $\deg(D) = 0$ to eliminate $n_\infty$ before building the
function; this is necessary, since the very first nontrivial principal divisor
$\divv(z) = (0)-(\infty)$ has $\infty$ in its support.
\end{remark}

\subsection{The main theorem}

\begin{theorem}\label{thm:main}
On $\PP$ one has $\Prin(\PP) = \ker(\deg)$. That is, a divisor is principal if
and only if its degree is zero.
\end{theorem}

\begin{proof}
Proposition~\ref{prop:degprin} gives $\Prin(\PP) \subseteq \ker(\deg)$, and
Proposition~\ref{prop:degzero} gives $\ker(\deg) \subseteq \Prin(\PP)$.
\end{proof}

\begin{smjcorollary}\label{cor:generators}
$\ker(\deg)$ is generated as a group by the differences $(p) - (q)$ with
$p, q \in \PP$, and each such difference is principal.
\end{smjcorollary}

\begin{proof}
The proof of Proposition~\ref{prop:degzero} exhibits an arbitrary degree-zero
divisor as an integer combination of the differences $(a_k) - (\infty)$, so these
already generate. Each $(p)-(q)$ has degree $0$ and is therefore principal by
Theorem~\ref{thm:main}; explicitly $(a)-(b) = \divv\big((z-a)/(z-b)\big)$ for
finite $a, b$, and $(a) - (\infty) = \divv(z-a)$.
\end{proof}

\begin{remark}[Why $\PP$ is special]\label{rem:genus}
Theorem~\ref{thm:main} says that on the Riemann sphere the degree is the
\emph{only} obstruction to a divisor being principal. This is a genuinely special
feature of $\PP$ and not a general phenomenon. On a compact Riemann surface $X$
of genus $g$, the degree-zero part of the class group is
\[
\Cl^0(X) \;\cong\; \operatorname{Jac}(X),
\]
a complex torus of dimension $g$ called the Jacobian of $X$. Thus a degree-zero
divisor is principal precisely when its class in this torus vanishes, which for
$g \geq 1$ is a nontrivial condition on a positive-dimensional family of classes.
The sphere is the case $g = 0$, where the torus is a single point and the
condition is vacuous --- which is exactly Theorem~\ref{thm:main}. Everything
below should be read as the genus-zero shadow of a much larger theory; see
Section~\ref{sec:future}.
\end{remark}

% ===============================================================
